\documentclass[10pt,a4paper]{amsart}
\usepackage[margin=19mm]{geometry}
\usepackage{amsmath,amssymb,mathtools}
\usepackage[scr=rsfs,cal=euler]{mathalfa}
\usepackage{xfrac}
\usepackage[unicode,hidelinks,bookmarks=false]{hyperref}
\newcommand{\ie}{i.e.\ }
\newcommand{\cf}{cf.\ }
\newcommand{\resp}{respectively\ }
\newcommand{\Subsection}{\subsection}
\providecommand{\nobreakdash}{\nobreak\hbox{-}}
\setlength{\emergencystretch}{2em}
\multlinegap0pt
\usepackage{amssymb}
\usepackage{dsfont}
\usepackage{xcolor}
\usepackage{caption}
\newtheorem{lemma}{Lemma}
\newtheorem{theorem}{Theorem}
\newcommand{\A}{{\mathbb{A}}}
\newcommand{\B}{{\mathbb{B}}}
\DeclareMathOperator{\Span}{Span}
\DeclareMathOperator{\Spin}{Spin}
\newcommand{\X}{{\mathbb{X}}}
\newcommand{\Z}{{\mathbb{Z}}}
\DeclareMathOperator{\coker}{coker}
\renewcommand{\frak}{\mathfrak}
\DeclareMathOperator{\im}{Im}
\DeclareMathOperator{\rank}{rank}
\newcommand{\wh}{\widehat}
\newcommand{\wt}{\widetilde}
\title{Heegaard Floer Surgery Formula and Cosmetic Surgeries}
\author{ALAN DU}
\date{Source: arXiv:2401.10395v2 (CC BY 4.0)}
\begin{document}
\maketitle

\noindent\emph{Source and licence.} Heegaard Floer Surgery Formula and Cosmetic Surgeries. Version arXiv:2401.10395v2 (CC BY 4.0), distributed by arXiv under the Creative Commons Attribution 4.0 International licence, http://creativecommons.org/licenses/by/4.0/, as recorded in the arXiv licence metadata for this exact version and shown on the abstract page https://arxiv.org/abs/2401.10395v2. Full licence notice: \texttt{licenses/arxiv-2401.10395-CC-BY-4.0.txt}.\par\smallskip

\noindent\textit{Editorial scope.} Counted unit: the article's theorem thm:rational-surgery-formula (source0.tex lines 442--445). The article's complete original proof of it is delivered with it (source0.tex lines 672--718, one \texttt{proof} environment of 47 lines, which the article places away from the statement). No mathematics is written, repaired or re-derived: the statement and the proof are the article's own lines, in the article's own notation, and no retained line is substituted or re-flowed.  Retained uncounted as the source-local context the proof needs: lines 246--248; lines 252--254; lines 628--630; lines 644--649.  Nothing was omitted from any retained span, so the package carries no omission record.  No retained line contains a citation, so the wrapper carries no bibliography and no bibliography entry.  No retained line contains a figure environment, an image-inclusion command or drawing code, so no image is delivered and the package declares no asset dependency.  Editorial formatting only, in addition: an AMS \texttt{amsart} preamble that restates the article's own declarations and macros used by the retained lines, namely the environments \texttt{lemma}, \texttt{theorem} on separate counters, together with the standard packages those lines need.  The retained environments print in this document as Lemma 1, Theorem 1 in order of appearance, whereas the article prints the same items with the numbering of its own full document.  The complete native source of the article as distributed in the arXiv e-print for this exact version is bundled unchanged as \texttt{sources/153184/source0.tex}.
\begin{lemma}\label{lem:v-h-equiv}
For all $s\in\Z$, $v_s^+$ and $h_s^+$ are chain homotopy equivalent to $h_{-s}^+\circ\phi$ and $v_{-s}^+\circ\phi$, respectively.
\end{lemma}

\begin{lemma}\label{lem:genus-finiteness}
Let $g$ be the genus of $K$, then for all $\delta\ge 0$ and $s\ge g+\delta$, $v_s^\delta$, and $h_{-s}^\delta$ induce isomorphisms on homology, and $v_{-s-1}^\delta$ and $h_{s+1}^\delta$ are nullhomotopic.
\end{lemma}

\begin{lemma}\label{lem:cancellation-general}
Let $K\subset Y$ be a blow-up bounding knot, $p,q>0$ relatively prime, and let $x\in H_*(\wh A_{j/q})$ for some $j$. If $\lfloor j/q\rfloor\ge 0$, there exists $\alpha\in H_*(\oplus_{i\ge j+p}\wh A_{i/q})$ such that $(\wh D_{p/q})_*(x+\alpha)=(\wh v_{j/q})_*(x)$ (where we think of $H_*(\wh B_{j/q})$ as including into $H_*(\wh\B)$). Similarly, if $\lfloor j/q\rfloor\le 0$, there exists $\alpha\in H_*(\oplus_{i\le j-p}\wh A_{i/q})$ such that $(\wh D_{p/q})_*(x+\alpha)=(\wh h_{j/q})_*(x)$.
\end{lemma}

\begin{align}\label{eq:rank-formula-pos-general}
\begin{split}
    \rank \wh{HF}(Y_{p/q}(K)) &= q(\rank\ker(\wh v_0)_*+b-\rank(\wh v_0)_*) \\
    &+2q\sum_{s=1}^{g-1}(\rank\ker(\wh v_s)_*+b-\rank(\wh v_s)_*)+ 2t_K^{p/q}-pb,
\end{split}
\end{align}

\begin{theorem}\label{thm:rational-surgery-formula}
Let $K\subset Y$ be a null-homologous knot, and let $p,q$ be a pair of relatively prime integers. Then, for each $\frak t\in\Spin^c(Y)$, $i\in\Z/p\Z$, there is an isomorphism of $\Z[U]$-modules
$$H_*(\mathbb X_{p/q,i,\frak t}^+)\cong HF^+(Y_{p/q}(K),i,\frak t).$$
\end{theorem}

\begin{proof}
By Theorem~\ref{thm:rational-surgery-formula}, $\wh{HF}(Y_{p/q}(K))\cong H_*(\wh{\mathbb X}_{p/q})$, the homology of the mapping cone of $\wh D_{p/q}$ from $\wh{\A}=\oplus_{i\in\Z/p\Z}\wh\A_i$ to $\wh{\B}=\oplus_{i\in\Z/p\Z}\wh\B_i$ (we have omitted the slope $p/q$ from the notation and summed over all $\Spin^c$ structures on $Y$). So
$$\wh\A=\bigoplus_{i=0}^{p-1}\prod_{r\in\Z}\wh A_{(i+pr)/q}=\prod_{j\in\Z}\wh A_{j/q}$$ has $q$ copies of $\wh A_s$ for each $s\in\Z$, and under the map $\wh v$, each $\wh A_{j/q}$ corresponds to a copy of $\wh B$ in $\wh\B$, which we shall denote $\wh B_{j/q}$. Then $\wh h_{j/q}$ maps $\wh A_{j/q}$ to $\wh B_{(j+p)/q}$. Thus, for each $s\in\Z$, there are $p$ copies of $\wh B_s$ in the codomain of $\wh h_{s-1}$ and $q-p$ copies of $\wh B_s$ in the codomain of $\wh h_s$.

The finiteness of $H_*(\wh\X_{p/q})$ is guaranteed by Lemma~\ref{lem:genus-finiteness}; in fact, it implies that $\wh\X_{p/q}$ is isomorphic to a truncated version of the mapping cone
$$\wh\X_{p/q}^c=MC(\wh D_{p/q}^c:\wh\A_{p/q}^c\to\wh\B_{p/q}^c),$$ for some large $c$, where
$$\wh\A_{p/q}^c=\bigoplus_{j=-qc+1}^{qc-1}\wh A_{j/q}=\bigoplus_{s=-c+1}^{c-1}\bigoplus_{i=0}^{q-1}\wh A_s, \hspace{5mm} \wh\B_{p/q}^c=\bigoplus_{j=-qc+p+1}^{qc-1}\wh B_{j/q},$$
and $\wh D_{p/q}^c$ is the restriction of $\wh D_{p/q}$ to $\wh A_{p/q}^c$, with the exception that for any $x\in\wh A_{j/q}$ with $-qc+1\le j\le-qc+p$, $\wh D_{p/q}^c(x)=\wh h_{j/q}(x)$, and for any $x\in\wh A_{j/q}$ with $qc-p\le j\le qc-1$, $\wh D_{p/q}^c(x)=\wh v_{j/q}(x)$. Note that it suffices to take $c\ge g+p/q+1$, and
$$\rank H_*(\wh\A_{p/q}^c) = q\sum_{s=-c+1}^{c-1}\rank H_*(\wh A_s), \hspace{5mm} \rank H_*(\wh\B_{p/q}^c) = qb(2c-1)-pb.$$

We construct a basis for $\ker(\wh D_{p/q})_*$ of the entire mapping cone, which has the same rank as $\ker(\wh D_{p/q}^c)_*$ for large $c$ by Lemma~\ref{lem:genus-finiteness}. For $j\ge 0$, let
$$\{ x_{j/q}^i\}_{i=1}^{\rank\ker(\wh v_{j/q})_*}$$
be a basis of $\ker(\wh v_{j/q})_*$. For each such $x_{j/q}^i$, choose an $\alpha_{j/q}^i\in H_*(\oplus_{k\ge j+p}\wh A_{k/q})$ according to Lemma~\ref{lem:cancellation-general} such that $\wt x_{j/q}^i=x_{j/q}^i+\alpha_{j/q}^i$ is in $\ker(\wh D_{p/q})_*$. For $j<0$, let 
$$\{x_{j/q}^i\}_{i=1}^{\rank\ker(\wh h_{j/q})_*}$$ 
be a basis of $\ker(\wh h_{j/q})_*$. For each such $x_{j/q}^i$, choose a $\beta_{j/q}^i\in H_*(\oplus_{k\le j-p}\wh A_{k/q})$ according to Lemma~\ref{lem:cancellation-general} such that $\wt x_{j/q}^i=x_{j/q}^i+\beta_{j/q}^i$ is in $\ker(\wh D_{p/q})_*$. Lastly, for $0\le j\le p-1$, let 
$$\{y_{j/q}^i\}_{i=1}^{\rank\left(\im(\wh v_{j/q})_* \cap \im(\wh h_{(j-p)/q})_*\right)}$$
be a basis of a subspace of $H_*(\wh A_{j/q})$ that is mapped isomorphically to $\im(\wh v_{j/q})_* \cap \im(\wh h_{(j-p)/q})_*$ under $(\wh v_{j/q})_*$. Choose $z_{j/q}^i\in H_*(\wh A_{(j-p)/q})$ such that $(\wh v_{j/q})_*(y_{j/q}^i)=(\wh h_{(j-p)/q})_*(z_{j/q}^i)$. Choose $\gamma_{j/q}^i\in H_*(\oplus_{k\ge j+p}\wh A_{k/q})$ and $\delta_{j/q}^i\in H_*(\oplus_{k\le j-p}\wh A_{k/q})$ according to Lemma~\ref{lem:cancellation-general} such that $\wt y_{j/q}^i=y_{j/q}^i+z_{j/q}^i+\gamma_{j/q}^i+\delta_{j/q}^i$ is in $\ker(\wh D_{p/q})_*$.

The set $$S=\{\wt x_{j/q}^i\}\cup\{\wt y_{j/q}^i\}$$ is linearly independent, and it is finite by Lemma~\ref{lem:genus-finiteness}. Moreover, it spans $\ker(\wh D_{p/q})_*$: let $z\in\ker(\wh D_{p/q})_*$ be nonzero, and suppose first that $z\in H_*(\oplus_{k\ge 0}\wh A_{k/q})$. Let $z_{j/q}$ denote the term of $z$ in $H_*(\wh A_{j/q})$, and let
$$j=\min\{k\ge 0\ \vert\ z_{k/q}\ne 0\}.$$ Then $z\in\ker(\wh D_{p/q})_*\cap H_*(\oplus_{k\ge j}\wh A_{k/q})$, so $z_{j/q}\in\Span \{x_{j/q}^i\}_{i=1}^{\rank\ker(\wh v_{j/q})_*}$: otherwise if $z_{j/q}\notin\ker(\wh v_{j/q})_*$, then since $z\in H_*(\oplus_{k\ge j}\wh A_{k/q})$, we have $z_{(j-p)/q}=0$ so $$(\wh v_{j/q})_*(z_{j/q})+(\wh h_{(j-p)/q})_*(z_{(j-p)/q})\ne 0\in H_*(\wh B_{j/q})$$ which means $z\notin\ker(\wh D_{p/q})_*$. Thus, we can subtract some combination of the basis elements $\wt x_{j/q}^i$ to eliminate the $H_*(\wh A_{j/q})$ component of $z$:
$$z-\sum_{i\in J}\wt x_{j/q}^i\in\ker(\wh D_{p/q})_*\cap H_*(\oplus_{k>j}\wh A_{k/q})$$ for some indexing set $J\subset\{1,\dots,\rank\ker(\wh v_{j/q})_*\}$. Inductively, we can then write 
$$z=\sum_{k\in I}\sum_{i_k\in J_k}\wt x_{k/q}^{i_k}$$ for some finite indexing sets $I\subset\{k\ge j\}$, $J_k\subset\{1,\dots,\rank\ker(\wh v_{k/q})_*\}$. Similarly, if $z\in H_*(\oplus_{k<0}\wh A_{k/q})$, then we can write
$$z=\sum_{k\in I}\sum_{i_k\in J_k}\wt x_{k/q}^{i_k}$$ for some finite indexing set $I\subset\{k\le j\}$, $J_k\subset\{1,\dots,\rank\ker(\wh h_{k/q})_*\}$.

Now for the general case $z\in\ker(\wh D_{p/q})_*$, we claim that for each $0\le j\le p-1$, $$z_{j/q}\in\Span\left(\{y_{j/q}^i\}\cup\{x_{j/q}^i\}\right).$$ Indeed, suppose this is not true for some $0\le j\le p-1$, then $(\wh v_{j/q})_*(z_{j/q})\notin\im(\wh h_{(j-p)/q})_*$, so $(\wh D_{p/q})_*(z)\ne 0$, a contradiction. Thus,
\begin{align*}
z-\sum_{j=0}^{p-1}\left(\sum_{i\in J_j}\wt y_{j/q}^i+\sum_{i\in J_j'}\wt x_{j/q}^i\right) &\in \left(\ker(\wh D_{p/q})_*\cap H_*(\oplus_{k\ge p}\wh A_{k/q})\right) \\
&\hspace*{5mm}+\left(\ker(\wh D_{p/q})_*\cap H_*(\oplus_{k<0}\wh A_{k/q})\right).
\end{align*} By the above argument, $z$ is in the span of $S$.

Therefore, $S$ is a basis for $\ker(\wh D_{p/q})_*$, so since $\rank\ker(\wh v_s)_*=\rank\ker(\wh h_{-s})_*$ by the hat version of Lemma~\ref{lem:v-h-equiv} and these are $0$ when $s\ge g$ by Lemma~\ref{lem:genus-finiteness},
$$\rank\ker(\wh D_{p/q})_*=q\rank\ker(\wh v_0)_*+2q\sum_{s=1}^{g-1}\rank\ker(\wh v_s)_*+t_K^{p/q}.$$

Since $\rank(\wh v_s)_*=\rank H_*(\wh A_s)-\rank\ker(\wh v_s)_*$ and $\rank\coker(\wh D_{p/q}^c)_*=\rank H_*(\wh\B_{p/q}^c)-\rank H_*(\wh\A_{p/q}^c)+\rank\ker(\wh D_{p/q}^c)_*$ for some large $c$, it follows that
\begin{align*}
    \rank\coker(\wh D_{p/q}^c)_* &= b(2c-1)q-pb-q\sum_{s=-c+1}^{c-1} \rank H_*(\wh A_s) \\
        &\hspace*{5mm}+ q\rank\ker(\wh v_0)_*+2q\sum_{s=1}^{g-1}\rank\ker(\wh v_s)_*+t_K^{p/q} \\
        &= q(b-\rank H_*(\wh A_0)+\rank\ker(\wh v_0)_*) \\
        &\hspace*{5mm}+ 2q\sum_{s=1}^{g-1}(b-\rank H_*(\wh A_s)+\rank\ker(\wh v_s)_*) +t_K^{p/q}-pb \\
        &= q(b-\rank(\wh v_0)_*)+2q\sum_{s=1}^{g-1}(b-\rank(\wh v_s)_*)+t_K^{p/q}-pb,
\end{align*}
where the second equality uses that $\rank H_*(\wh A_s)=b$ for all $|s|\ge g$ by Lemmas~\ref{lem:v-h-equiv} and ~\ref{lem:genus-finiteness}. Therefore, because 
$$\rank(H_*(\wh\X_{p/q}))=\rank\ker(\wh D_{p/q})_*+\rank\coker(\wh D_{p/q})_*$$ and $$\rank\ker(\wh D_{p/q})_*=\rank\ker(\wh D_{p/q}^c)_*,\hspace*{5mm}\rank\coker(\wh D_{p/q})_*=\rank\coker(\wh D_{p/q}^c)_*$$ for large enough $c$, equation~\ref{eq:rank-formula-pos-general} holds.



\end{proof}

\end{document}
